\section{Results}
\label{sec:results}

\hypertarget{feature-philosophy-and-model-selection-k1-k2-debutanizer}{%
\subsection{Feature philosophy and model selection (K1, K2 --- Debutanizer)}\label{feature-philosophy-and-model-selection-k1-k2-debutanizer}}

Table~\ref{tab:t1} / Figure~\ref{fig:f2} report the feature ablation under the identical blocked
time-series CV protocol (5 folds, transport lag 15, plain linear regression, scaler fit
per fold-train). The four arms differ only in where model complexity comes from.

\textbf{T1 --- Debutanizer feature ablation (blocked CV; held-out test n = 360).}

\begin{table}[t]\centering\small
\caption{Debutanizer feature ablation under the identical blocked-CV protocol (held-out test $n=360$).}\label{tab:t1}
\begin{tabular}{@{}
  >{\raggedright\arraybackslash}p{(\columnwidth - 8\tabcolsep) * \real{0.2000}}
  >{\raggedright\arraybackslash}p{(\columnwidth - 8\tabcolsep) * \real{0.2000}}
  >{\raggedright\arraybackslash}p{(\columnwidth - 8\tabcolsep) * \real{0.2000}}
  >{\raggedright\arraybackslash}p{(\columnwidth - 8\tabcolsep) * \real{0.2000}}
  >{\raggedright\arraybackslash}p{(\columnwidth - 8\tabcolsep) * \real{0.2000}}@{}}
\toprule\noalign{}
\begin{minipage}[b]{\linewidth}\raggedright
arm
\end{minipage} & \begin{minipage}[b]{\linewidth}\raggedright
k
\end{minipage} & \begin{minipage}[b]{\linewidth}\raggedright
blocked-CV R²
\end{minipage} & \begin{minipage}[b]{\linewidth}\raggedright
worst fold
\end{minipage} & \begin{minipage}[b]{\linewidth}\raggedright
test R²
\end{minipage} \\
\midrule\noalign{}

\bottomrule\noalign{}
\endlastfoot
u5 only (backbone) & 1 & +0.034 ± 0.438 & −1.697 & +0.395 \\
physics only & 3 & +0.120 ± 0.479 & −1.770 & +0.440 \\
physics + u5 (deployed) & 4 & \textbf{+0.145 ± 0.419} & −1.494 & \textbf{+0.476} \\
lagged raw (kitchen sink) & 126 & −1.604 ± 1.421 & −6.248 & +0.221 \\
\end{tabular}
\end{table}

Three observations. First, the thermodynamically grounded set (bubble-point estimate,
relative volatility α(T), stripping factor) dominates on every metric while using 4
features instead of 126: complexity set by physics beats complexity set by validation.
Second, the kitchen sink's failure is not a cross-validation artifact --- its held-out
test R² (+0.221) is less than half the deployed set's (+0.476), so the overfitting is
visible out-of-sample as well. Third, the negative folds in \emph{every} arm are the
cross-regime selection lottery that motivates the protocol: under naive adjacent
validation these regime shifts are averaged away and selection becomes seed-dependent;
under blocked CV they surface as the per-fold spread that the one-standard-error rule
then acts on. The deployed arm's row is bit-identical to the static row of the
independent bias-update evaluation (Section 5.2), an internal consistency check across
two separately written evaluation scripts.

The transport-lag diagnosis precedes all of this: at lag 0 the physics bridge appears
to fail (static R² ≈ 0.02), which a naive reading would call a physics error; lag
scanning recovers u5 at lag 15 as a stable interior maximum (r² ≈ 0.51). The lag-zero
failure is diagnostic of transport delay, not of the thermodynamics (results ledger, \S 1A).

\hypertarget{drift-handling-under-delayed-labels-k2-debutanizer-tep}{%
\subsection{Drift handling under delayed labels (K2 --- Debutanizer, TEP)}\label{drift-handling-under-delayed-labels-k2-debutanizer-tep}}

Figure~\ref{fig:f4} shows the held-out residual trace before and after the open-loop bias-update
\citep{shardt2016} at the documented analyzer delay (λ = 0.1 CV-selected on the pool, test
untouched; θ = 4). The static model's per-fold R² spread of {[}+0.224, \textbf{−1.494},
+0.574, +0.728, +0.695{]} (CV +0.145 ± 0.419) collapses to {[}+0.614, +0.487, +0.684,
+0.704, +0.751{]} (CV +0.648 ± \textbf{0.046}) --- a 9× reduction in cross-fold standard error,
which is the robustness claim: the update rescues the regime-shifted fold (−1.49 →
+0.49) without sacrificing the good ones. On the held-out test window the same
mechanism moves R² from +0.476 to \textbf{+0.857} (Δ +0.381). The correction also exceeds
the best \emph{constant} per-fold offset (oracle CV +0.540), confirming that within-fold
drift tracking --- not merely offset removal --- is doing work.

The same recipe applied unchanged to the TEP regimes reproduces the signature
(blocked-CV −0.089 ± 0.493 static; bias-update recovery to CV +0.19 at θ = 5, +0.45 at
θ = 2) --- the methodology transfers across process topology before any model migration
is attempted (results ledger, \S 1C).

\hypertarget{calibrated-uncertainty-under-drift-and-delay-k3-tep-synthetic}{%
\subsection{Calibrated uncertainty under drift and delay (K3 --- TEP, synthetic)}\label{calibrated-uncertainty-under-drift-and-delay-k3-tep-synthetic}}

Static split conformal \citep{papadopoulos2002} assumes exchangeability; process drift
breaks it. On a controlled synthetic stress test (residual scale ×3 mid-stream, target
0.90), split coverage collapses to 0.403 while ACI \citep{gibbs2021aci} holds 0.899
and EnbPI \citep{xu2021enbpi} recovers with FIFO lag (0.832) --- Figure~\ref{fig:f6}. The real TEP
regimes exhibit the production version of this failure: not a collapse but
\emph{inconsistency} (Table~\ref{tab:t2} / Figure~\ref{fig:f7}).

\textbf{T2 --- Empirical coverage (target 0.90), TEP regimes × delay θ; n\_test = 300/regime.}

\begin{table}[t]\centering\small
\caption{Empirical coverage against the 0.90 target on real TEP regimes, by construction and label delay $\theta$ ($n_\mathrm{test}=300$ per regime). Width in parentheses.}\label{tab:t2}
\begin{tabular}{@{}llll@{}}
\toprule\noalign{}
regime & construction & θ = 2 cov (width) & θ = 5 cov (width) \\
\midrule\noalign{}

\bottomrule\noalign{}
\endlastfoot
mode 1 & raw + split & 0.847 (1.79) & 0.847 (1.79) \\
& corrected + ACI & \textbf{0.900} (1.72) & 0.897 (2.04) \\
& EnbPI & 0.897 (2.11) & 0.897 (2.11) \\
mode 2 & raw + split & 0.847 (2.11) & 0.847 (2.11) \\
& corrected + ACI & \textbf{0.900} (1.90) & 0.903 (2.18) \\
& EnbPI & 0.857 (2.26) & 0.857 (2.26) \\
mode 3 & raw + split & 0.957 (3.18) & 0.957 (3.18) \\
& corrected + ACI & \textbf{0.900} (2.01) & 0.900 (2.48) \\
& EnbPI & 0.873 (2.49) & 0.873 (2.49) \\
\end{tabular}
\end{table}

The static baseline swings from under-coverage (0.847, modes 1--2) to over-coverage
(0.957, mode 3) --- invalid in both directions, and the over-covering regime pays for its
nominal-looking number with the widest intervals (3.18). The corrected + ACI
construction is regime-uniform at the target --- coverage range {[}0.897, 0.903{]} across all
six regime × delay cells --- with widths 1.7--2.5, and is robust to the label delay (θ
enters through the bias-update; ACI is driven by the 1D.1b protocol, γ = 0.05). EnbPI
sits between the two, structurally limited by its FIFO residual window on these
mild-drift regimes (its batch-size parameter is flat here: coverage 0.897--0.910 over
s ∈ \{1, 10, 25, 50\}).

\hypertarget{model-migration-across-operating-regimes-k4-tep}{%
\subsection{Model migration across operating regimes (K4 --- TEP)}\label{model-migration-across-operating-regimes-k4-tep}}

Migration is evaluated as within-process regime migration (mode 1 → modes 2, 3) under
the two-layer composition --- offline migration plus the Section 5.2 online bias-update
(λ = 0.3, θ = 2) --- with n = 8 random-subset repeats per data fraction; literal
cross-process parameter migration is mathematically inapplicable here because the
source and target input spaces are disjoint (ADR-009). The metric is \textbf{data
efficiency}: the ratio of data fractions at which from-scratch and migrated models
reach 90\% of the from-scratch ceiling. (The stricter ceiling-crossover statistic
flipped between 20\% and 100\% across seeds in preliminary runs and is not headlined.)

Figure~\ref{fig:f5}: Yan functional scale-bias correction \citep{yan2011} --- a per-input affine plus a
zero-mean GP bias on the source predictor --- reaches 90\% of the ceiling at \textbf{5\%} of
target data versus \textbf{50\%} from scratch on \emph{both} targets: \textbf{10.0× data efficiency},
with calibrated 95\% GP posterior intervals (coverage 92--98\%) that are appropriately
wider for the larger mode-3 regime gap (widths 2.84--3.96 vs 2.65--3.16). Luo matrix
scale-bias correction \citep{luo2015} yields exactly 1.0× on both targets: for a linear
source its per-input scaling absorbs the source weights and the method is from-scratch
regression in disguise --- the migrated and from-scratch curves coincide to ±0.001--0.006
at every fraction, the empirical face of an identity we also verify analytically.
Output-only correction \citep{lu2008similarity} lands between (1.0× / 4.0×). A bare-migration control
(no bias-update layer) caps every migrated curve below the 90\% threshold on mode 2:
the two-layer composition is load-bearing, because within-regime drift otherwise
swamps the cross-regime correction.

\hypertarget{the-negative-control-no-physics-anchors-k5-secom}{%
\subsection{The negative control: no physics anchors (K5 --- SECOM)}\label{the-negative-control-no-physics-anchors-k5-secom}}

SECOM inverts every assumption the framework exploits: 590 anonymized features over
1,567 samples (p ≈ n), 4.5\% missing cells, no thermodynamic structure. Framed as
virtual metrology (target x59, selected by stated criteria --- missingness ≤ 5\%, maximal
\textbar point-biserial r\textbar{} = 0.156 with pass/fail; audit table in the run output), the
elastic-net path under blocked CV explodes at weak shrinkage --- CV R² reaching −7 × 10⁴
with standard errors of the same magnitude (Figure~\ref{fig:f3}), the p ≈ n amplification of the
Section 5.1 lottery. The one-SE rule selects maximal shrinkage (α* = 10, 6/441 nonzero
coefficients) and contains the damage; the resulting sensor still fails as a point
predictor: test R² = \textbf{−1.84} raw, −0.16 after the bias-update absorbs a large
calibration offset --- even corrected, it cannot beat predicting the mean.

Coverage, however, survives (Table~\ref{tab:t3}): corrected + ACI delivers 0.910 / 0.915 (θ =
2/5) against the 0.90 target at mean width 12.8, where the static split over-covers
(0.953) at width 20.2 --- 37\% wider. \textbf{Conformal validity is model-agnostic; point
accuracy is what the physics anchors buy.} On a process where the sensor cannot be
made accurate, the framework still tells the truth about its own uncertainty. The
fail-enrichment of conformal misses is inconclusive (9 failing lots in the test
window) and is reported only as such.

\textbf{T3 --- Cross-dataset summary (the negative-control row is the point).}

\begin{table}[t]\centering\small
\caption{Cross-dataset summary; the negative-control row is the point.}\label{tab:t3}
\begin{tabular}{@{}
  >{\raggedright\arraybackslash}p{(\columnwidth - 6\tabcolsep) * \real{0.2500}}
  >{\raggedright\arraybackslash}p{(\columnwidth - 6\tabcolsep) * \real{0.2500}}
  >{\raggedright\arraybackslash}p{(\columnwidth - 6\tabcolsep) * \real{0.2500}}
  >{\raggedright\arraybackslash}p{(\columnwidth - 6\tabcolsep) * \real{0.2500}}@{}}
\toprule\noalign{}
\begin{minipage}[b]{\linewidth}\raggedright
dataset
\end{minipage} & \begin{minipage}[b]{\linewidth}\raggedright
physics anchors
\end{minipage} & \begin{minipage}[b]{\linewidth}\raggedright
point R² (raw → corrected)
\end{minipage} & \begin{minipage}[b]{\linewidth}\raggedright
ACI coverage (target 0.90)
\end{minipage} \\
\midrule\noalign{}

\bottomrule\noalign{}
\endlastfoot
Debutanizer & yes (VLE bridge) & +0.476 → \textbf{+0.857} & --- (Section 5.3 protocol on TEP) \\
TEP regimes & yes (process features) & +0.31 → CV +0.45 (θ=2) & \textbf{0.897--0.903}, all cells \\
SECOM & \textbf{none} & −1.84 → −0.16 & \textbf{0.910 / 0.915} \\
\end{tabular}
\end{table}

\hypertarget{deployment-evidence-k5}{%
\subsection{Deployment evidence (K5)}\label{deployment-evidence-k5}}

The serving stack (Figure~\ref{fig:f1}b) was exercised end-to-end: 243 tests pass (CI: lint,
unit, Docker build, container smoke on \texttt{/health} and \texttt{/predict}), and client-observed
latency on commodity hardware is p50 1.23 ms / p90 1.51 / p99 1.97 over 10⁴ single-row
predictions --- including HTTP and serialization --- with batch-32 amortizing to ≈46 µs
per row and the delayed-label \texttt{/label} round-trip at p50 0.97 ms. The 200 ms
soft-sensor budget is met with two orders of magnitude of headroom; the path is
I/O-bound by construction, since the deployed object is a linear model, an EWMA, and a
quantile lookup. Correctness under delayed labels is enforced structurally: the
coverage indicator for a late label is computed against the interval \emph{stored when that
sample was predicted}, not the interval current at arrival time (Figure~\ref{fig:f1}b, dashed) ---
by which time α\_t and the score window have moved on.

% --- figures first referenced in this section ---

\begin{figure}[t]
\centering
\includegraphics[width=0.6\textwidth]{F2_feature_ablation}
\caption{\textbf{Feature ablation on the debutanizer (identical blocked-CV protocol).} Bars: blocked-CV R² mean ± SE; red dashes: worst fold; k = number of features. The physics-anchored set with the raw lagged temperature (k = 4) dominates (CV +0.145 ± 0.419, test +0.476), while the 126-feature lagged kitchen sink collapses (CV −1.604 ± 1.421, worst fold −6.248) and its failure is visible out-of-sample as well (test +0.221): complexity set by physics beats complexity set by validation.}
\label{fig:f2}
\end{figure}

\begin{figure}[t]
\centering
\includegraphics[width=0.6\textwidth]{F3_secom_cv_path}
\caption{\textbf{SECOM elastic-net CV path (p ≈ n selection lottery).} Mean blocked-CV −R² (log scale) ± SE along the regularization path, complexity increasing to the right. The path explodes by four orders of magnitude (to −R² ≈ 7 × 10⁴) at weak shrinkage; the one-standard-error rule (blue) selects maximal shrinkage α* = 10 (6/441 nonzero coefficients), adjacent to the best CV mean (red), containing the lottery that pure CV-mean selection would re-enter.}
\label{fig:f3}
\end{figure}

\begin{figure}[t]
\centering
\includegraphics[width=\textwidth]{F4_bias_trace}
\caption{\textbf{Held-out residual trace, raw vs.~bias-corrected (debutanizer).} Residuals on the 360-sample test window before (red) and after (blue) the open-loop EWMA bias update (λ = 0.1 CV-selected on the pool; θ = 4). The correction removes the calibration offset and tracks within-window drift, lifting held-out R² from +0.476 to +0.857; the residual sd drops accordingly. The same mechanism reduces cross-fold SE 9× in CV (Table~\ref{tab:t1} vs.~§5.2).}
\label{fig:f4}
\end{figure}

\begin{figure}[t]
\centering
\includegraphics[width=\textwidth]{F5_data_efficiency}
\caption{\textbf{Data efficiency of regime migration (TEP, mode 1 → modes 2 and 3;
two-layer = migration + online bias update, n = 8 repeats).} Held-out R² vs.~fraction of target-regime data for the migrated sensor (blue, ± sd) and from-scratch training (red); dashed line: 90\% of the from-scratch ceiling. Yan functional scale-bias migration reaches the threshold at 5\% of target data vs.~50\% from scratch on both targets --- 10.0× data efficiency --- with calibrated GP posterior intervals (coverage 92--98\%) that widen appropriately for the larger mode-3 regime gap.}
\label{fig:f5}
\end{figure}

\begin{figure}[t]
\centering
\includegraphics[width=0.6\textwidth]{F6_synthetic_stress}
\caption{\textbf{Coverage under a controlled drift stress (synthetic; target 0.90).} Residual scale is tripled mid-stream. Static split conformal collapses to 0.403 coverage; adaptive conformal inference re-attains 0.899 by adapting the working miscoverage level; EnbPI recovers to 0.832 with the lag of its FIFO residual window. The controlled version of the real-regime inconsistency in Figure~\ref{fig:f7}.}
\label{fig:f6}
\end{figure}

\begin{figure}[t]
\centering
\includegraphics[width=\textwidth]{F7_tep_coverage}
\caption{\textbf{Empirical coverage on real TEP regimes (target 0.90, dashed; n = 300 per
regime; θ ∈ \{2, 5\}).} Static split conformal on raw residuals is invalid in both directions --- under-covering modes 1--2 (0.847) and over-covering mode 3 (0.957, at the widest intervals, 3.18). The corrected + ACI construction is regime-uniform at the target (range {[}0.897, 0.903{]} across all six regime × delay cells) at widths 1.7--2.5; EnbPI sits between. Validity here is a property of the composition: bias correction first, conformal calibration on what remains.}
\label{fig:f7}
\end{figure}
